% =====================================================================
%  Projeto FEMM --- Conversão da Energia (UFMG)
%  Parte analítica (itens a e b): fatores de enrolamento e espectro
%  de B_g(theta) idealizado.
%
%  Bruno --- Eng. Elétrica
%  Prof. Braz J. Cardoso Filho
% =====================================================================

\documentclass[12pt,a4paper]{article}

% ---------- Pacotes ----------
\usepackage[utf8]{inputenc}
\usepackage[T1]{fontenc}
\usepackage[brazil]{babel}
\usepackage[margin=2.2cm]{geometry}
\usepackage{amsmath,amssymb}
\usepackage{booktabs}
\usepackage{array}
\usepackage{enumitem}
\usepackage{xcolor}
\usepackage{hyperref}
\usepackage{titlesec}
\usepackage{parskip}
\usepackage{siunitx}
\usepackage{caption}
\usepackage{pgfplots}
\pgfplotsset{compat=1.18}
\usepgfplotslibrary{groupplots}
\usetikzlibrary{patterns}

\hypersetup{
    colorlinks=true,
    linkcolor=blue!60!black,
    urlcolor=blue!60!black,
    citecolor=blue!60!black
}

\titleformat{\section}{\large\bfseries\color{blue!50!black}}{\thesection}{1em}{}
\titleformat{\subsection}{\normalsize\bfseries}{\thesubsection}{1em}{}

\setlist[itemize]{leftmargin=*,itemsep=2pt,topsep=2pt}
\setlist[enumerate]{leftmargin=*,itemsep=2pt,topsep=2pt}

\sisetup{output-decimal-marker={,}, group-separator={\,}, group-minimum-digits=4}

% Comando auxiliar para ressaltar identidades importantes
\newcommand{\eqimp}[1]{\boxed{\,#1\,}}

% =====================================================================
\begin{document}
% =====================================================================

\title{\vspace{-1.5cm}\textbf{Projeto computacional --- Análise harmônica de uma máquina de indução trifásica}\\[0.2em]
\large Parte~I: análise analítica (itens \textbf{a} e \textbf{b})}
\author{Bruno Ribeiro e Arnaldo --- UFMG, Eng. Elétrica\\
\small Prof. Braz J. Cardoso Filho --- Conversão da Energia}
\date{Versão preliminar --- \today}

\maketitle

% ---------------------------------------------------------------------
\section*{Resumo}

Este documento desenvolve a parte analítica do projeto, cujo objetivo é estudar o espectro espacial da força magnetomotriz (FMM) e da densidade de fluxo no entreferro $B_g(\theta)$ para três configurações de enrolamento de estator de uma máquina trifásica de indução: (i) camada simples com passo pleno, (ii) camada dupla com passo pleno, (iii) camada dupla com passo encurtado $\beta=5/6$. Toda a análise se apoia nas equações (12a)--(12c) do formulário (Apêndice~B do Chapman) e no Capítulo~4 do Fitzgerald. A modelagem em FEMM (itens c e d) será feita em uma etapa posterior, a partir das mesmas configurações aqui caracterizadas.

% ---------------------------------------------------------------------
\section{Dados do projeto}

\begin{table}[h]
\centering
\caption{Parâmetros geométricos e elétricos extraídos do enunciado.}
\renewcommand{\arraystretch}{1.1}
\begin{tabular}{ll}
\toprule
\textbf{Parâmetro} & \textbf{Valor} \\
\midrule
Número de polos $P$              & 6 \\
Número de fases $m$              & 3 \\
Bobinas por polo por fase $q$    & 4 \\
Número de ranhuras $Q = m\,P\,q$ & \textbf{72} \\
Tensão terminal                  & \SI{127}{V_{rms}} \\
Frequência elétrica $f_e$        & \SI{60}{Hz} \\
Entreferro $g$                   & \SI{0.5}{mm} \\
Comprimento do pacote $L$        & \SI{140}{mm} \\
Densidade de fluxo no núcleo $B_{\text{máx}}$ & \SI{1.7}{T} \\
$B_g$ de pico (assumido)         & $\approx \SI{0.9}{T}$ \\
Aço elétrico                     & E250 (M250-50A), \SI{0.5}{mm} \\
\midrule
Passo polar $\tau_p$ (em ranhuras)        & $Q/P = 12$ \\
Ângulo elétrico entre ranhuras $\alpha$   & $\dfrac{P}{2}\cdot\dfrac{360^{\circ}}{Q} = 15^{\circ}$ \\
Largura do cinturão (\emph{phase belt}) $q\alpha$ & $60^{\circ}$ \\
Velocidade síncrona (eq.~4a) $n_s = \tfrac{120 f_e}{P}$ & $\SI{1200}{RPM}$ \\
\bottomrule
\end{tabular}
\end{table}

\noindent A figura indicativa do enunciado mostrava 36 ranhuras, mas os parâmetros do texto fixam $Q=72$ ranhuras --- adota-se $Q=72$ como valor de referência.

% ---------------------------------------------------------------------
\section{Configurações de enrolamento analisadas}

A análise comparativa cobre as três configurações abaixo. Em todas, o passo polar é de \mbox{12 ranhuras}; o que muda é o número de camadas e o passo da bobina.

\begin{enumerate}
    \item[\textbf{(I)}] \textbf{Camada simples, passo pleno.}
    Cada ranhura aloja uma bobina inteira; o passo da bobina é igual ao passo polar ($\beta=12/12=1$, ou \mbox{$180^{\circ}$ elétricos}).
    
    \item[\textbf{(II)}] \textbf{Camada dupla, passo pleno.}
    Mesma característica eletromagnética da (I) na escala da fundamental, porém com cada ranhura recebendo dois lados de bobina (camadas superior e inferior). Vantagem construtiva: padronização das bobinas pré-moldadas.
    
    \item[\textbf{(III)}] \textbf{Camada dupla, passo encurtado $\beta=5/6$.}
    O passo da bobina é reduzido a $5\tau_p/6 = 10$ ranhuras (ou $150^{\circ}$ elétricos). Esta é a escolha clássica em motores trifásicos para \emph{atacar simultaneamente} os harmônicos de quinta e sétima ordem, que são os mais expressivos depois da fundamental \cite{chapman2013}.
\end{enumerate}

\noindent\textbf{Observação importante.} Configurações (I) e (II) compartilham os mesmos fatores $k_d$, $k_p$ e portanto $k_w$ --- a passagem de camada simples para camada dupla é uma decisão construtiva, não eletromagnética (na escala da fundamental espacial idealizada). Em todas as tabelas adiante, as colunas (I) e (II) são apresentadas separadamente apenas para fidelidade ao escopo do projeto, mas ambas se baseiam em $\beta=1$.

% ---------------------------------------------------------------------
\section{Item (a) --- Fatores $k_d$, $k_p$ e $k_w$}

\subsection{Equações de partida}

A formulação clássica para os fatores de enrolamento de bobinas distribuídas em ranhuras igualmente espaçadas, generalizada para uma harmônica espacial de ordem $\nu$, segue diretamente das equações~(12a)--(12c) do formulário, aplicadas a cada harmônico ímpar:
\begin{align}
k_{d,\nu} &= \frac{\sin\!\bigl(\nu\,q\,\alpha/2\bigr)}{q\,\sin\!\bigl(\nu\,\alpha/2\bigr)}, \label{eq:kd}\\[2pt]
k_{p,\nu} &= \sin\!\left(\nu\,\beta\,\frac{\pi}{2}\right), \label{eq:kp}\\[2pt]
k_{w,\nu} &= k_{p,\nu}\cdot k_{d,\nu}, \label{eq:kw}
\end{align}
em que $\alpha=15^{\circ}$ é o passo de ranhura em graus elétricos, $q=4$ é o número de ranhuras por polo por fase e $\beta$ é o passo relativo da bobina (1 para passo pleno; $5/6$ para o encurtamento clássico). O sinal de $k_{p,\nu}$ é mantido para preservar a fase relativa de cada harmônico; o módulo $|k_{w,\nu}|$ é o que determina a amplitude do harmônico na onda de FMM.

\subsection{Cálculo da fundamental ($\nu=1$)}

\paragraph{Fator de distribuição.} Substituindo $q=4$, $\alpha=15^{\circ}$ na eq.~\eqref{eq:kd} para $\nu=1$,
\[
k_{d,1} = \frac{\sin(4\cdot 15^{\circ}/2)}{4\,\sin(15^{\circ}/2)}
       = \frac{\sin(30^{\circ})}{4\sin(7{,}5^{\circ})}
       = \frac{0{,}5000}{4\cdot 0{,}1305}
       \approx \eqimp{0{,}9577}.
\]

\paragraph{Fator de passo.} Para passo pleno, $k_{p,1}=\sin(90^{\circ})=1$. Para $\beta=5/6$,
\[
k_{p,1} = \sin\!\left(\frac{5}{6}\cdot 90^{\circ}\right) = \sin(75^{\circ}) \approx \eqimp{0{,}9659}.
\]

\paragraph{Fator de enrolamento total.} Pela eq.~\eqref{eq:kw}:
\[
k_{w,1}^{\beta=1}\!=\!0{,}9577,\qquad
k_{w,1}^{\beta=5/6}\!=\!0{,}9659\cdot 0{,}9577 \approx \eqimp{0{,}9250}.
\]

\noindent A redução é de apenas \(\sim 3{,}4\%\) na fundamental --- preço pequeno pelos benefícios em harmônicos.

\subsection{Tabela completa (até $\nu=13$)}

A Tabela~\ref{tab:fatores} reúne os valores de $k_{d,\nu}$, $k_{p,\nu}$ e $k_{w,\nu}$ para os harmônicos ímpares relevantes. As triplens ($\nu=3,9,15,\ldots$) aparecem por completude --- elas \emph{não circulam} na FMM trifásica equilibrada (Seção~\ref{sec:itemb}), independentemente do enrolamento.

\begin{table}[h]
\centering
\caption{Fatores de enrolamento para $q=4$, $\alpha=15^{\circ}$.}
\label{tab:fatores}
\renewcommand{\arraystretch}{1.15}
\begin{tabular}{c c c c c c}
\toprule
$\nu$ & $k_{d,\nu}$ & $k_{p,\nu}\ (\beta=1)$ & $k_{p,\nu}\ (\beta=5/6)$ & $k_{w,\nu}\ (\beta=1)$ & $k_{w,\nu}\ (\beta=5/6)$\\
\midrule
1  & $+0{,}9577$ & $+1$ & $+0{,}9659$ & $+0{,}9577$ & $+0{,}9250$ \\
3  & $+0{,}6533$ & $-1$ & $-0{,}7071$ & $-0{,}6533$ & $-0{,}4619$ \\
5  & $+0{,}2053$ & $+1$ & $+0{,}2588$ & $+0{,}2053$ & $+0{,}0531$ \\
7  & $-0{,}1576$ & $-1$ & $+0{,}2588$ & $+0{,}1576$ & $-0{,}0408$ \\
9  & $-0{,}2706$ & $+1$ & $-0{,}7071$ & $-0{,}2706$ & $+0{,}1913$ \\
11 & $-0{,}1261$ & $-1$ & $+0{,}9659$ & $+0{,}1261$ & $-0{,}1218$ \\
13 & $+0{,}1261$ & $+1$ & $-0{,}9659$ & $+0{,}1261$ & $-0{,}1218$ \\
\bottomrule
\end{tabular}
\end{table}

\subsection{Discussão dos resultados}

\begin{itemize}
    \item \textbf{Fundamental.} A passagem para $\beta=5/6$ custa apenas $\sim 3{,}4\%$ no fator de enrolamento, valor sempre considerado aceitável.
    
    \item \textbf{Quinta e sétima harmônicas.} O encurtamento $5/6$ é \emph{otimizado} para esses harmônicos: $k_{p,5}=k_{p,7}=\sin(75^{\circ}\cdot 5)=\sin(75^{\circ}\cdot 7)=\sin(15^{\circ})\approx 0{,}259$. O resultado é uma queda de cerca de \textbf{74\%} no fator de passo desses harmônicos, e proporcionalmente em $k_w$:
    \[
    \frac{|k_{w,5}^{5/6}|}{|k_{w,5}^{1}|} \approx \frac{0{,}053}{0{,}205} \approx 0{,}258,
    \quad
    \frac{|k_{w,7}^{5/6}|}{|k_{w,7}^{1}|} \approx \frac{0{,}041}{0{,}158} \approx 0{,}259.
    \]
    
    \item \textbf{Terceira e nona (triplens).} O fator $k_{w}$ aparece grande, mas o cancelamento na soma trifásica torna esses harmônicos eletricamente irrelevantes na FMM resultante (vide Seção~\ref{sec:trif}).
    
    \item \textbf{Décima primeira e décima terceira.} Não são afetadas pelo encurtamento $5/6$ na mesma intensidade --- $|k_{p,11}|=|k_{p,13}|=\sin(75^{\circ})=0{,}9659$ (idêntico ao da fundamental). É um conhecido efeito colateral da escolha $5/6$: a janela ``deixa passar'' os harmônicos $11$ e $13$ por simetria. Outras escolhas (como $11/12$) atacam esses pares mas pioram outros.
\end{itemize}

% ---------------------------------------------------------------------
\section{Item (b) --- Espectro analítico de \texorpdfstring{$B_g(\theta)$}{Bg(theta)}}
\label{sec:itemb}

\subsection{Da FMM de uma fase à FMM trifásica girante}

A onda de FMM produzida por uma única fase é estacionária no espaço e pulsa no tempo. Para a fase~a, em uma máquina de $P$ polos, escreve-se (Eq.~4.6 do Fitzgerald, generalizada para harmônicos):
\begin{equation}
F_a(\theta_{ae},t) = \sum_{\nu \,\text{ímpar}} F_{\nu}\,\cos(\nu\,\theta_{ae})\cos(\omega_e t),
\label{eq:Fa}
\end{equation}
com amplitude por harmônico
\begin{equation}
F_{\nu} = \frac{4}{\pi}\,\frac{k_{w,\nu}\,N_{\text{fase}}}{\nu\,P}\,I_{\max}.
\label{eq:Fnu}
\end{equation}
A dependência $\propto 1/\nu$ é o decaimento natural da série de Fourier de uma onda quase-retangular gerada por bobinas concentradas, atenuado adicionalmente pelo fator $k_{w,\nu}$ que captura a distribuição em $q$ ranhuras e o eventual encurtamento de passo.

\subsection{Composição trifásica e cancelamento das triplens}
\label{sec:trif}

Aplicando a identidade $\cos a\cos b = \tfrac{1}{2}[\cos(a-b)+\cos(a+b)]$, cada harmônico de cada fase decompõe-se em duas ondas progressivas (uma direta e uma reversa). Somando as contribuições das três fases (deslocadas de $\pm 120^{\circ}$ no espaço e no tempo), obtém-se o resultado clássico (Seção~4.5.2 do Fitzgerald):
\begin{equation}
F_{\text{trif}}(\theta_{ae},t) = \frac{3}{2}\sum_{\nu \in \mathcal{H}^+} F_{\nu}\,\cos(\nu\theta_{ae}-\omega_e t)
                                + \frac{3}{2}\sum_{\nu \in \mathcal{H}^-} F_{\nu}\,\cos(\nu\theta_{ae}+\omega_e t),
\end{equation}
em que
\begin{itemize}
    \item $\mathcal{H}^+ = \{1,7,13,19,\ldots\}$ -- harmônicos que giram \emph{no mesmo sentido} da fundamental;
    \item $\mathcal{H}^- = \{5,11,17,23,\ldots\}$ -- harmônicos que giram \emph{em sentido oposto};
    \item $\mathcal{H}^0 = \{3,9,15,\ldots\}$ -- as triplens, que \emph{somam zero} nas três fases.
\end{itemize}

\noindent O fato relevante é, portanto, que \textbf{a FMM trifásica equilibrada não contém harmônicos de ordem múltipla de três}, mesmo quando $k_{w,3}$ é grande em uma fase isolada. A partir daqui, toda a análise espectral ignora as triplens.

\subsection{Densidade de fluxo no entreferro}

Sob a hipótese clássica de $\mu_{\text{ferro}} \to \infty$ e entreferro uniforme de espessura $g$, a aplicação da Lei de Ampère a um contorno que cruza o entreferro duas vezes resulta em
\begin{equation}
H_g(\theta) = \frac{F(\theta)}{g}, \qquad B_g(\theta) = \mu_0 H_g(\theta) = \frac{\mu_0\,F(\theta)}{g}.
\end{equation}
Logo, $B_g$ tem o \emph{mesmo} espectro espacial da FMM, escalado por $\mu_0/g$. A amplitude do $\nu$-ésimo harmônico de $B_g$, normalizada pela fundamental, é simplesmente
\begin{equation}
\eqimp{\;\frac{|B_{g,\nu}|}{|B_{g,1}|} \;=\; \frac{|k_{w,\nu}|}{\nu\,|k_{w,1}|}\;,\qquad \nu \notin \mathcal{H}^0.}
\label{eq:espectro}
\end{equation}

\subsection{Aplicação numérica}

Adotando $|B_{g,1}| = \SI{0.9}{T}$ como amplitude de referência, a Tabela~\ref{tab:espectro} apresenta o espectro analítico para as três configurações.

\begin{table}[h]
\centering
\caption{Espectro analítico de $B_g(\theta)$, em tesla, para as três configurações ($|B_{g,1}|=\SI{0.9}{T}$ por hipótese). Triplens omitidas (canceladas pela soma trifásica).}
\label{tab:espectro}
\renewcommand{\arraystretch}{1.15}
\begin{tabular}{c c c c c c}
\toprule
$\nu$ & $|k_{w,\nu}|/(\nu|k_{w,1}|)\ (\beta=1)$ & $|B_{g,\nu}|\ (\beta=1)$ &
        $|k_{w,\nu}|/(\nu|k_{w,1}|)\ (\beta=5/6)$ & $|B_{g,\nu}|\ (\beta=5/6)$ \\
\midrule
1  & $1{,}000$ & $0{,}9000$ & $1{,}000$ & $0{,}9000$ \\
5  & $0{,}0429$ & $0{,}0386$ & $0{,}0115$ & $0{,}0103$ \\
7  & $0{,}0235$ & $0{,}0212$ & $0{,}0063$ & $0{,}0057$ \\
11 & $0{,}0120$ & $0{,}0108$ & $0{,}0120$ & $0{,}0108$ \\
13 & $0{,}0101$ & $0{,}0091$ & $0{,}0101$ & $0{,}0091$ \\
\bottomrule
\end{tabular}
\end{table}

\noindent Resultados-chave:
\begin{itemize}
    \item A quinta harmônica passa de $\SI{38.6}{mT}$ (passo pleno) para $\SI{10.3}{mT}$ ($\beta=5/6$): \textbf{atenuação de $\approx 73\%$}.
    \item A sétima harmônica passa de $\SI{21.2}{mT}$ para $\SI{5.7}{mT}$: \textbf{atenuação de $\approx 73\%$}.
    \item As décima primeira e décima terceira harmônicas \emph{permanecem inalteradas} entre as duas configurações --- consequência direta de $|k_{p,11}|=|k_{p,13}|=|k_{p,1}|$ para $\beta=5/6$.
\end{itemize}

\subsection{Visualização do espectro}

A Figura~\ref{fig:espectro} mostra o espectro normalizado das três configurações em barras agrupadas. A escala vertical é logarítmica para tornar visíveis os harmônicos de alta ordem.

\begin{figure}[h]
\centering
\begin{tikzpicture}
\begin{axis}[
    width=14cm, height=8cm,
    ybar=2pt,
    bar width=8pt,
    enlarge x limits=0.12,
    ymode=log,
    log basis y=10,
    ymin=0.005, ymax=2,
    ylabel={$|B_{g,\nu}|/|B_{g,1}|$ \ (escala log)},
    xlabel={Ordem do harmônico espacial $\nu$},
    symbolic x coords={1,5,7,11,13},
    xtick=data,
    ymajorgrids=true,
    grid style={dashed,gray!40},
    legend style={at={(0.98,0.98)}, anchor=north east, font=\small, draw=gray!50, fill=white, fill opacity=0.92, text opacity=1},
    legend cell align={left},
    title={Espectro analítico normalizado de $B_g(\theta)$ --- triplens omitidas}
]
\addplot+[fill=blue!30, draw=blue!60!black] coordinates {
    (1, 1.0000) (5, 0.0429) (7, 0.0235) (11, 0.0120) (13, 0.0101)
};
\addplot+[fill=orange!40, draw=orange!70!black, postaction={pattern=north east lines, pattern color=orange!90!black}] coordinates {
    (1, 1.0000) (5, 0.0429) (7, 0.0235) (11, 0.0120) (13, 0.0101)
};
\addplot+[fill=red!40, draw=red!70!black] coordinates {
    (1, 1.0000) (5, 0.0115) (7, 0.0063) (11, 0.0120) (13, 0.0101)
};
\legend{{(I) Cam.\ simples \,$\beta=1$}, {(II) Cam.\ dupla \,$\beta=1$}, {(III) Cam.\ dupla \,$\beta=5/6$}}
\end{axis}
\end{tikzpicture}
\caption{Amplitudes relativas dos harmônicos espaciais de $B_g$ no entreferro (escala logarítmica). As configurações (I) e (II) coincidem analiticamente (a hachura na (II) é apenas para diferenciá-la visualmente); a configuração (III) suprime fortemente $\nu=5$ e $\nu=7$ mas preserva $\nu=11$ e $\nu=13$.}
\label{fig:espectro}
\end{figure}

% ---------------------------------------------------------------------
\section{Antecipação dos harmônicos de ranhura (item d)}

Embora o item (d) só possa ser fechado após a FFT do FEMM, é instrutivo já registrar as ordens esperadas dos harmônicos de ranhura, que são fenômenos puramente geométricos --- não dependem do passo da bobina nem do número de camadas \cite[\S B.5]{chapman2013}.

\paragraph{Fórmula.} As frequências espaciais introduzidas pela presença de $Q$ ranhuras uniformemente distribuídas em uma máquina de $P$ polos são
\begin{equation}
\nu_{\text{ranhura}} = \frac{2\,Q\,M}{P} \pm 1, \qquad M = 1, 2, 3, \ldots
\label{eq:slot}
\end{equation}

\paragraph{Aplicação numérica.} Para $Q=72$, $P=6$ ($p=P/2=3$):
\[
\nu_{\text{ranhura}} = \frac{72\,M}{3} \pm 1 = 24M \pm 1
\;\Rightarrow\;
\boxed{\;\nu = 23,\,25\;\text{(}M=1\text{)};\; \nu = 47,\,49\;\text{(}M=2\text{)};\; \ldots\;}
\]

\noindent Na frequência elétrica de \SI{60}{Hz}, esses harmônicos correspondem a $1380$ e $\SI{1500}{Hz}$ na tensão induzida \cite[Eq.~B-26]{chapman2013}. Eles devem aparecer claramente na FFT de $B_g$ extraída do FEMM e \textbf{não} são suprimidos pela escolha de passo --- exigem outras estratégias (inclinação dos condutores do rotor ou enrolamento fracionário).

% ---------------------------------------------------------------------
\section{Síntese e próximas etapas}

\paragraph{O que foi estabelecido.}
\begin{itemize}
    \item Fatores $k_d$, $k_p$ e $k_w$ calculados para harmônicos $\nu = 1$ a $13$, nas três configurações de enrolamento.
    \item Fundamental: $k_{w,1}=0{,}9577$ (passo pleno) e $0{,}9250$ (encurtado $5/6$).
    \item Espectro analítico de $B_g$ derivado da FMM trifásica girante: $|B_{g,\nu}|/|B_{g,1}| = |k_{w,\nu}|/(\nu|k_{w,1}|)$, com triplens canceladas.
    \item Atenuação de $\approx 73\%$ em $\nu=5$ e $\nu=7$ pelo encurtamento; $\nu=11$ e $13$ não atenuados.
    \item Harmônicos de ranhura previstos em $\nu = 23, 25, 47, 49, \ldots$ ($M=1,2,\ldots$).
\end{itemize}

\paragraph{O que falta (FEMM).}
\begin{enumerate}
    \item Modelagem da geometria com $Q=72$, entreferro $g=\SI{0.5}{mm}$ e curva $BH$ não-linear do aço E250.
    \item Configuração de excitação no instante de pico da fase~A: $i_A=I_p$, $i_B=i_C=-I_p/2$. Calibração de $N\!\cdot\! I$ para que $|B_{g,1}|\approx \SI{0.9}{T}$.
    \item Extração de $B_g(\theta)$ ao longo de um arco no centro do entreferro.
    \item FFT espacial e comparação direta com a Tabela~\ref{tab:espectro} e a Figura~\ref{fig:espectro}, com atenção ao surgimento dos harmônicos de ranhura previstos pela eq.~\eqref{eq:slot}.
\end{enumerate}

\paragraph{Pontos a confirmar com o Prof. Braz.}
\begin{itemize}
    \item Hipótese de trabalho do passo encurtado: $\beta=5/6$ (\mbox{passo $=10$ ranhuras}). Alternativas razoáveis: $11/12$ ou $9/12$.
    \item Valor de $|B_{g,1}|$ a ser empregado na calibração da fonte FMM no FEMM.
\end{itemize}

% ---------------------------------------------------------------------
\begin{thebibliography}{9}
\bibitem{fitzgerald2014} A.~E.~Fitzgerald, C.~Kingsley Jr., S.~D.~Umans. \textit{Máquinas Elétricas}, 7\textsuperscript{a} ed. AMGH, 2014. (Capítulo 4: ondas girantes de FMM.)
\bibitem{chapman2013} S.~J.~Chapman. \textit{Fundamentos de Máquinas Elétricas}, 5\textsuperscript{a} ed. AMGH, 2013. (Apêndice~B: fatores de enrolamento e harmônicos de ranhura.)
\bibitem{formulario} Formulário P2 --- Conversão da Energia (Prof. Braz J. Cardoso Filho), eqs.~(12a)--(12c).
\end{thebibliography}

\end{document}